\chapter{Implementation}\label{chap:implementation}
\renewcommand{\proves}{\vdash}

This section describes a concrete implementation of a restricted subset of VAMP's, with the goal of preserving the core dynamics while making simulation and learning tractable and realistic.

\subsection{Assumptions}

We implement a family of environments under the following simplifying assumptions.

\paragraph{Finite Formula Universe.}
Each environment instance is defined over a finite universe of formulas (nodes) $\mathcal{F}$.
All sets derived from $\mathcal{F}$ (ghost/concrete partitions, candidate actions, etc.)
are therefore finite.

\paragraph{Finite Agent Universe.}
We assume $\mathcal{N} = \{1,\cdots,n\}$ for some fixed $n \in \mathbb N$.

% \paragraph{Agent Capabilities.}
% We assume agents $\alpha \in \mathcal{N}$ are equipped with stationary proving and conjecturing abilities. Namely, although agents may learn strategies and update their internal query models based on experience, we freeze their proving and conjecturing kernels: $\mathcal{K}_\mathrm{proof}^\alpha$ $\mathcal{N} = \{1,\cdots,n\}$ for some fixed $n \in \mathbb N$.

\paragraph{Truth and Difficulty Initialization.}

We assume we may initialize the VAMP with a latent truth labeling over formulas $T:\mathcal F\to\{0,1\}$, such that for all $\phi\in\mathcal F$, $T(\phi)+T(\neg\phi)=1$.

We interpret $T(\phi)=1$ as ``$\phi$ is true/valid in the instance'' and $T(\phi)=0$ as ``$\phi$ is false''. For concision, we define $\mathcal{F}_\top := \{\phi \in \mathcal F : T(\phi)=1\}$. Of course, if $[\Gamma \proves \phi]\in\mathcal S$ is a resolved sequent then $\phi \in \mathcal{F}_\top$.
This truth labeling is fixed over time and unobserved by agents except through interaction outcomes.

Additionally, we initialize a static latent difficulty map
\[
d:\mathcal F_\top\to (0,1)
\]
used by the proof kernel as a (unobservable) hyperparameter in assigning success/failure probabilities heterogeneously. $d$ is supported over $\{\phi \in \mathcal F : T(\phi)=1\}$.




\paragraph{Static dependency structure and metadata.}
We assume there exists a fixed, instance-dependent \emph{ground truth} dependency mapping over formulas
\[
\delta^\star : \mathcal{F}_\top\to \mathcal{P}_{\mathrm{fin}}(\mathcal{F})
\]

where we overload this notation to note for resolved sequents $[\Gamma \proves \phi]$ via $\delta^\star([\Gamma \proves \phi]) := \delta^\star(\phi) \subseteq \Gamma$ and such that the induced dependency digraph is acyclic. Then we may define for $s \in \mathcal{RS}_{t_0}^\alpha$ for all $t \ge t_0$, $\delta^\alpha_t(s) = \delta^\star(s)$.

We note this is simply an environment modeling choice. We do not claim mathematical proofs are unique; rather, it yields clean initialization and unambiguous progress signals.

\paragraph{Utility Graph.}
Each instance of a VAMP initializes from a weighted directed graph
\[
G_U := (\mathcal F_\top, E_U, w),\qquad w:E_U\to(0,1],
\]
where $w(\psi,\phi)$ measures the ``utility'' of $\psi$ toward resolving or reasoning about $\phi$. This is necessarily an informal notion that simply serves to allow us to highly related concepts that are not formally dependent on one another. We extend $w$ by $w(\psi,\phi)=0$ when $(\psi, \phi)\notin E_U$.

We additionally note that:
\[w(\psi, \phi)=1 \iff \psi \in \delta^*(\phi)\]
(So, if $E^\star = \{(\psi,\phi) \in E_U : w(\psi,\phi)=1\} \subseteq E_U$, then the subgraph $G^\star = (\mathcal{F}_\top,E^\star)$ is a DAG).


\paragraph{Sequent Specification.}
To prevent an exponential blowup of action space size due to agents' selection of sequents $[\Gamma \proves \phi]$ for $\Gamma \subseteq_{\mathrm{fin}} \mathcal{F}$, we automatically construct antecedent context based on the local library state. Namely, our implemented interface restricts agents to selecting only a consequent formula $\phi$. The antecedent used to evaluate the attempt is then deterministically constructed from the agent's current local library state.

For agent $\alpha$ at time $t$, define the set of locally available resolved formulas
\[
\mathcal{R}^\alpha_t
:= \Big\{\psi\in\mathcal{F}:\exists \Gamma'\subseteq_{\mathrm{fin}}\mathcal{F}\text{ such that }[\Gamma'\proves \psi]\in\mathcal{RS}^\alpha_t\Big\},
\]

Now, we may redefine the (non-market) action space:
\begin{align*}
\widetilde{\mathcal{A}_{\mathrm{nm}}^\alpha}
:=
&
\{\mathrm{Prove}\}\times(\mathcal{F}\times\mathbb{N}^+)
\\&\sqcup
\{\mathrm{Conj}\}\times(\mathcal{F}\times\mathbb{N}^+)
\\&\sqcup
\{\mathrm{Pub}\}\times\mathcal{F}
\\&\sqcup
\{\mathrm{Qry}\}\times \mathcal{S}
\\&\sqcup\{\bot_{\mathrm{nm}}\}
\end{align*}

We note that the query action still relies on sequents rather than formulas, though we will shortly bound the admissible input sequents for this action in particular to allow for polynomial growth of the admissible action space size, while maintaining the ability of agents to ``suppose lemma $\psi$ is proven'' when running a query. W fix a hyperparameter $h_{\mathrm{max}}$ to denote this boundedness of the admissible antecedents' extension beyond $\mathcal{R}^\alpha_t$.

Additionally, as the market mechanism action space is abstractly defined, we note that it is up to the user specification to decide whether similar implementation changes must be changed in $\mathcal{A}_{\mathrm{mkt}}$ as well.

Additionally, we update the specification of the admissibility criteria, where for $w_t \in \mathcal{W}$ and $\alpha \in \mathcal{N}$, if $J^\alpha_t=\bot$, then:
\begin{itemize}
\item For $a^\alpha_t = (\mathrm{Prove},(\phi,\tau))$ and $a^\alpha_t = (\mathrm{Conj},(\phi,\tau))$:
\[
\widetilde{\mathrm{Adm}^\alpha}(w_t,a^\alpha_t)=1 \iff \phi \in \mathcal{C}_t^\alpha
\]
\item For $a^\alpha_t = (\mathrm{Pub},\phi)$: \[\widetilde{\mathrm{Adm}^\alpha}(w_t,a^\alpha_t)=1 \iff \phi \in \mathcal{C}_t^\alpha\]
\item For $a^\alpha_t = (\mathrm{Qry},[\Gamma \proves \phi])$: 
\begin{align*}
    \widetilde{\mathrm{Adm}^\alpha}(w_t,a_t^\alpha)=1 \iff &[\Gamma \proves \phi] \in \mathcal{CS}_t^\alpha \\&\land (|\Gamma \setminus \mathcal{R}^\alpha_t| \le h_{\mathrm{max}})
\end{align*}
\end{itemize}
Where in all other cases, $\widetilde{\mathrm{Adm}^\alpha}(w_t,a_t^\alpha) =\mathrm{Adm}^\alpha(w_t,a_t^\alpha)$. 

Then, given $w_t \in \mathcal{W}$ and admissible $\widetilde{a_t^\alpha} \in \widetilde{\mathcal{A}^\alpha_t}(w_t)$, we may recover a corresponding $a_t^\alpha \in \mathcal{A}^\alpha_t(w_t)$, defined as follows:
\begin{itemize}
    \item If $\widetilde{a_t^\alpha} = (\mathrm{Prove},\phi,\tau)$ then $a_t^\alpha = (\mathrm{Prove},[\mathcal{R}_t^\alpha \proves \phi],\tau)$.
    \item If $\widetilde{a_t^\alpha} = (\mathrm{Conj},\phi,\tau)$ then $a_t^\alpha = (\mathrm{Conj},[\mathcal{R}_t^\alpha \proves \phi],\tau)$.
    \item If $\widetilde{a_t^\alpha} = (\mathrm{Pub},\phi)$ then $a_t^\alpha = (\mathrm{Pub},\{\phi\},\{[\mathcal{R}_t^\alpha \proves \phi]\}\cap \mathcal{R}^\alpha_t)$.
    \item Otherwise, $\widetilde{a_t^\alpha} = a_t^\alpha$.
\end{itemize}


We note that this notion of fixing a canonical local context $\mathcal{R}_t^\alpha$ effectively defines a bijection from (concrete) formulas to (concrete) sequents, in that if we define
\[
s^\alpha_t(\phi)\ :=\ [\Gamma^\alpha_t \proves \phi]\ \in\ \mathcal{S}.
\]
\[
\widetilde{\mathcal{CS}}^\alpha_t
:=\{\,s^\alpha_t(\phi): \phi\in\mathcal{C}^\alpha_t\,\},
\]
and the map $\phi\mapsto s^\alpha_t(\phi)$ is a bijection between $\mathcal{C}^\alpha_t$ and
$\widetilde{\mathcal{CS}}^\alpha_t$. %\textbf{TODO FORMALLY PROVE, AS THIS CS IS DIFFERENT THAN USUAL CS}

Additionally, we note that this assumption automatically implies that all resolved sequents are actually fully resolved.%: \textbf{TODO PROVE!!!}

% We additionally implement the proof kernel with a \emph{resolved-dependency gate}:
% a proof attempt on $\phi$ may only succeed if its returned dependency set $d$ satisfies
% \[
% d\subseteq \Gamma^\alpha_t=\mathrm{ResForm}^\alpha_t.
% \]
% Consequently, any newly resolved sequent $[\Gamma^\alpha_t\proves\phi]$ depends only on formulas that were
% already resolved at time $t$, and therefore (by induction on resolution time, or by acyclicity/topological
% ordering of the dependency DAG) every resolved sequent is automatically \emph{fully} resolved.
% Hence, in the implemented regime we may identify
% \[
% \mathcal{RS}^\alpha_t=\mathcal{FRS}^\alpha_t
% \qquad\text{and similarly for the public library } \mathcal{RS}^0_t=\mathcal{FRS}^0_t.
% \]


\subsection{Proof Kernel ($\mathcal{K}_{\text{proof}}$).}
The proof kernel models the outcome of spending compute budget on an admissible proof attempt. For an agent $\alpha$ and its library state $\mathcal{L}\in\mathscr{L}$, target sequent $s\in\mathcal S$ (though in practice $s \in \mathcal{CS}^{\alpha}$), and compute budget $\tau\in\mathbb N^+$, the kernel outputs a binary success indicator and set of dependencies. However, as we assume a global dependency set, we may simply consider $\mathcal{K}_{\text{proof}}$ outputting in $\{0,1\}$, with the dependency component being computed deterministically as $\delta^\star(s)$.

With that in mind, we now proceed to explicitly define our particular implementation of $\mathcal{K}_{\text{proof}}$, stepping through its parameterization and desired behaviors. Namely, we construct $\mathcal{K}_{\text{proof}}$ to have the following properties:

\paragraph{Bernoulli proof outcomes.}
We assume proof attempts have only binary stochasticity. Concretely, $\mathcal K^\alpha_{\mathrm{proof}}(\cdot\mid \mathcal{L}, s,\tau)$ is supported on exactly two atoms:
\[
\mathcal K^\alpha_{\mathrm{proof}}\big((0,\bot)\mid \mathcal{L}, s,\tau\big)=1-p^\alpha_{\mathrm{proof}}(s,\tau)\]
\[\mathcal K^\alpha_{\mathrm{proof}}\big((1,\delta^\star(s))\mid \mathcal{L}, s,\tau\big)=p^\alpha_{\mathrm{proof}}(s,\tau)
\]
and assigns zero mass to all $(1,d)$ with $d\neq \delta^\star(s)$.

\paragraph{Truth + Dependency Gate.} If $\phi$ is false (i.e. $T(\phi)=0$), the probability of success is $0$. Moreover, if $\delta^\star(\phi)\not\subseteq \mathcal{R}_t^\alpha$, success probability is 0. In other words, one cannot prove a theorem before its hard prerequisites.

Towards this, we may define a feasibility gate
\[
g([\Gamma\proves \phi])=\mathbf{1}[T(\phi)=1]\cdot \mathbf{1}[\delta^\star(\phi)\subseteq \Gamma],
\]

\paragraph{Budget Monotonicity and Diminishing Returns.} We enforce that $p^\alpha_{\mathrm{proof}}(s,\cdot)$ is nondecreasing, intuitively modeling that more cumulative compute can’t hurt success likelihood. Additionally, we hope to model diminishing returns with increasing budget such that marginal gains decrease with more budget. This serves to prevent undesirable strategies such as ``always dump infinite compute into one node''. A natural implementation of this is to model:
\[
p_{\text{proof}}^\alpha([\Gamma \proves \phi],\tau)=K_0\left(1-e^{-K_1\cdot h(\tau)}\right)
\]
For some $K_0,K_1\ge 0$, and a concave $h$. In practice, one often chooses $h(\tau) = \tau^\kappa$ governed by a hyperparameter $\kappa \in (0,1]$.

\paragraph{Difficulty Modulation.} We recall that we initialize $\mathcal{F}_\top$ with a latent difficulty map $d: \mathcal{F}_\top \to (0,1)$. Given the same time budget $\tau$, harder theorems are less likely to be solved. Moreover, difficulty should scale the rate (hazard) rather than just shifting probabilities ad hoc. Towards this notion, define:
\[
r_{\mathrm{diff}}(\phi)=e^{-\lambda_{\mathrm{diff}}\,d(\phi)}
\]
Governed by a difficulty-to-rate transform $\lambda_{\mathrm{diff}}>0$.

\paragraph{Utility Modulation.} We also recall that we initialize a utility graph $G_U$ over $\mathcal{F}_\top$ as a weighted digraph. Although the feasibility gate $g$ filters sequents $[\Gamma \proves \phi]$ based on the inclusion of the hard dependencies $\delta^\star$, we aim to model the notion that even once hard dependencies are met, having additional relevant resolved lemmas in $\Gamma$ should increase success probabilities. This too should be monotone and saturating. Towards this, we define:
\[
s_{\text{mass}}([\Gamma \proves \phi])=\sum_{\psi\in \Gamma} w(\psi,\phi)^\alpha
\]
with $\alpha\ge 1$, and a saturating transform
\[
s([\Gamma \proves \phi])=\log\big(1+s_{\text{mass}}([\Gamma \proves \phi])\big)
\]
This design preserves the distinction between having one useful theorem and many useful theorems, while preventing unbalanced growth in success rate. It is governed by a hyperparameter $\alpha \ge 1$ to affects the sensitivity to high-utility neighbors.

\paragraph{Heterogeneous Agent Capability Modulation.} We initialize agents with capability parameters:
\begin{itemize}
    \item $\rho_0^\alpha\ge 0$ representing agent $\alpha$'s baseline solve-rate
	\item $\rho_1^\alpha\ge 0$ representing agent $\alpha$'s amplification from high-utility support
\end{itemize}
With this in mind, a natural model of the effective proof rate of agent $\alpha$ on a particular target sequent is:
\[
\rho^\alpha([\Gamma \proves \phi])=r_\mathrm{diff}(\phi)\Big(\rho_0^\alpha + \rho_1^\alpha s([\Gamma \proves \phi])\Big)
\]

\subsubsection{Definition}

As such, we may explicitly implement:
\[
\mathcal K^\alpha_{\mathrm{proof}}\big((0,\bot)\mid \mathcal{L}, s,\tau\big)=1-p^\alpha_{\mathrm{proof}}(s,\tau)\]
\[\mathcal K^\alpha_{\mathrm{proof}}\big((1,\delta^\star(s))\mid \mathcal{L}, s,\tau\big)=p^\alpha_{\mathrm{proof}}(s,\tau)
\]
and assigns zero mass to all $(1,d)$ with $d\neq \delta^\star(s)$, where:
\[
p_{\text{proof}}^\alpha([\Gamma \proves \phi],\tau)=g([\Gamma \proves \phi])\left(1-e^{-\rho^\alpha([\Gamma \proves \phi])\cdot h(\tau)}\right)
\]
We note that $p_{\text{proof}}^\alpha(\cdot)\in[0,1]$ for all inputs, and success probability is strictly increasing and approaches $1$ as $\tau_{\mathrm{eff}}\to\infty$ on feasible inputs that pass $g$.

\subsection{Conjecture Kernel ($\mathcal{K}_{\text{conj}}$).}
The conjecture kernel models the outcome of allocating compute budget toward proposing a new formula. Given an agent $\alpha\in\mathcal N$, a library state $\mathcal L=(\mathcal C,\mathcal{RS},\delta,\Theta)\in\mathscr L$, an anchor sequent $s=[\Gamma\proves\phi]\in\mathcal S$ (though in practice $s\in\mathcal{CS}^\alpha$), and budget $\tau\in\mathbb N^+$, the conjecture kernel outputs either failure or success together with a proposed formula $\psi\in\mathcal F$. However, successful outputs are required to land in the current ghost set $\mathcal G(\mathcal{L}) :=\mathcal F\setminus \mathcal C$.

We now define our particular implementation of $\mathcal K^\alpha_{\mathrm{conj}}$, stepping through its desired behaviors and parameterization.

\paragraph{Bernoulli success with conditional proposal.}
We assume conjecturing has binary stochasticity at the success/failure level, followed by a proposal distribution on success. Concretely, there exists a success probability
\[
p^\alpha_{\mathrm{conj}}(\mathcal L,s,\tau)\in[0,1]
\]
and a conditional proposal distribution
\[
\pi^\alpha_{\mathrm{conj}}(\cdot\mid \mathcal L,s,\tau)
\]
supported on $\mathcal G(\mathcal L)$ such that
\[
\mathcal K^\alpha_{\mathrm{conj}}\big((0,\bot)\mid \mathcal L,s,\tau\big)
=
1-p^\alpha_{\mathrm{conj}}(\mathcal L,s,\tau),
\]
and for $\psi\in\mathcal F$,
\[
\mathcal K^\alpha_{\mathrm{conj}}\big((1,\psi)\mid \mathcal L,s,\tau\big)
=
p^\alpha_{\mathrm{conj}}(\mathcal L,s,\tau)\,\pi^\alpha_{\mathrm{conj}}(\psi\mid \mathcal L,s,\tau).
\]

\paragraph{Budget monotonicity and diminishing returns.}
As with the proof kernel, we require $p^\alpha_{\mathrm{conj}}(\mathcal L,s,\cdot)$ to be nondecreasing with diminishing returns. We therefore model
\[
p^\alpha_{\mathrm{conj}}(\mathcal L,s,\tau)
=
1-\exp\!\big(-\eta^\alpha(\mathcal L,s)\,h(\tau)\big),
\]
where $h:\mathbb N^+\to\mathbb R_{\ge 0}$ is increasing and concave. In practice, we again use
\[
h(\tau)=\tau^\kappa,\qquad \kappa\in(0,1].
\]

\paragraph{Heterogeneous agent conjecturing capability.}
We equip each agent $\alpha$ with conjecturing capability parameters
\[
\eta_0^\alpha,\eta_1^\alpha\ge 0,
\]
representing baseline conjecture rate and amplification from local graph opportunity, respectively.

\paragraph{Utility-guided relevance.}
We recall the instance utility graph
\[
G_U=(\mathcal F_\top,E_U,w),
\]
where $w(\psi,\phi)$ measures the utility of $\psi$ toward resolving or reasoning about $\phi$, extended by $0$ off $E_U$.

Intuitively, useful conjectures should depend on the role of the anchor consequent $\phi$ in the current local library. If $\phi$ is already resolved, then good conjectures should tend to extend outward from $\phi$; if $\phi$ is not yet resolved, then good conjectures should tend to provide lemmas that support $\phi$.

% Accordingly, define the local resolved formula set induced by $\mathcal{RS}$:
% \[
% \mathcal R(\mathcal L):=\Big\{\psi\in\mathcal F:\exists \Gamma'\subseteq_{\mathrm{fin}}\mathcal F,\; [\Gamma'\proves\psi]\in\mathcal{RS}\Big\},
% \]
and define the anchor-conditioned utility score
\[
q(\psi\mid \mathcal L,[\Gamma\proves\phi])
=
\mathbf 1[[\Gamma\proves\phi]\in\mathcal {RS}^\alpha]\,w(\phi,\psi)
+
\mathbf 1[[\Gamma\proves\phi]\notin\mathcal {RS}^\alpha]\,w(\psi,\phi).
\]
Thus, when $\phi$ is resolved, the score favors frontier extensions; when $\phi$ is unresolved, it favors potentially useful supporting lemmas.

\paragraph{Local opportunity modulation.}
We also want conjecturing to be easier in regions where there exist many promising unseen formulas. To capture this, define the local opportunity mass
\[
m^\alpha(\mathcal L,s)
:=
\sum_{\psi\in\mathcal G(\mathcal L)} \big(q(\psi\mid \mathcal L,s)\big)^{\beta_{\mathrm{conj}}},
\qquad
\beta_{\mathrm{conj}}\ge 1,
\]
and the saturating transform
\[
u^\alpha(\mathcal L,s):=\log\big(1+m^\alpha(\mathcal L,s)\big).
\]
We then define the effective conjecturing rate by
\[
\eta^\alpha(\mathcal L,s)
=
\eta_0^\alpha+\eta_1^\alpha\,u^\alpha(\mathcal L,s).
\]

\paragraph{Budget-improved proposal quality.}
Beyond increasing the chance of any successful conjecture, additional compute should also improve the quality of the proposed conjecture conditional on success. We model this by allowing the proposal distribution to become more concentrated around high-utility candidates as budget increases.

To do so, define a budget-dependent selectivity parameter
\[
\kappa^\alpha_{\mathrm{conj}}(\tau)
=
\kappa^\alpha_0+\kappa^\alpha_1 h(\tau),
\qquad
\kappa^\alpha_0,\kappa^\alpha_1\ge 0.
\]
Small values of $\kappa^\alpha_{\mathrm{conj}}(\tau)$ yield diffuse proposal distributions, while large values concentrate mass on high-score candidates.

\paragraph{Proposal distribution over all formulas, restricted to ghosts by renormalization.}
It is convenient to first define a raw proposal score over all formulas $\psi\in\mathcal F$, then restrict to the current ghost set induced by $\mathcal L$.

Let
\[
r^\alpha(\psi\mid \mathcal L,s,\tau)
:=
\exp\!\Big(\kappa^\alpha_{\mathrm{conj}}(\tau)\,\varphi(q(\psi\mid \mathcal L,s))\Big),
\]
where $\varphi:[0,1]\to\mathbb R_{\ge0}$ is a fixed monotone transform (for instance $\varphi(x)=x$ or $\varphi(x)=\log(1+x)$).

We then define the actual proposal distribution by renormalizing over the current ghost set:
\[
\pi^\alpha_{\mathrm{conj}}(\psi\mid \mathcal L,s,\tau)
=
\frac{r^\alpha(\psi\mid \mathcal L,s,\tau)\,\mathbf 1[\psi\in\mathcal G(\mathcal L)]}
{\sum_{\psi'\in\mathcal G(\mathcal L)} r^\alpha(\psi'\mid \mathcal L,s,\tau)}.
\]
Hence the conjecturer is specified globally over $\mathcal F$, but at runtime only samples from formulas that are currently ghost relative to $\mathcal L$.

\subsubsection{Definition}

Putting the above together, for an admissible conjecture attempt on $s=[\Gamma\proves\phi]$ with budget $\tau$, we define
% \[
% p^\alpha_{\mathrm{conj}}(\mathcal L,s,\tau)
% =
% 1-\exp\!\big(-\eta^\alpha(\mathcal L,s)\,h(\tau)\big),
% \]
% with
% \[
% \eta^\alpha(\mathcal L,s)
% =
% \eta_0^\alpha+\eta_1^\alpha\,\log\!\Big(1+\sum_{\psi\in\mathcal G(\mathcal L)}(q(\psi\mid \mathcal L,s))^{\beta_{\mathrm{conj}}}\Big),
% \]
% and
% \[
% \pi^\alpha_{\mathrm{conj}}(\psi\mid \mathcal L,s,\tau)
% =
% \frac{\exp\!\Big(\kappa^\alpha_{\mathrm{conj}}(\tau)\,\varphi(q(\psi\mid \mathcal L,s))\Big)\,\mathbf 1[\psi\in\mathcal G(\mathcal L)]}
% {\sum_{\psi'\in\mathcal G(\mathcal L)}
% \exp\!\Big(\kappa^\alpha_{\mathrm{conj}}(\tau)\,\varphi(q(\psi'\mid \mathcal L,s))\Big)}.
% \]

% Accordingly,
\[
\mathcal K^\alpha_{\mathrm{conj}}\big((0,\bot)\mid \mathcal L,s,\tau\big)
=
1-p^\alpha_{\mathrm{conj}}(\mathcal L,s,\tau),
\]
and for $\psi\in\mathcal F$,
\[
\mathcal K^\alpha_{\mathrm{conj}}\big((1,\psi)\mid \mathcal L,s,\tau\big)
=
p^\alpha_{\mathrm{conj}}(\mathcal L,s,\tau)\,\pi^\alpha_{\mathrm{conj}}(\psi\mid \mathcal L,s,\tau).
\]
By construction, $\pi^\alpha_{\mathrm{conj}}(\cdot\mid \mathcal L,s,\tau)$ is supported on $\mathcal G(\mathcal L)$, $p^\alpha_{\mathrm{conj}}(\mathcal L,s,\cdot)$ is monotone with diminishing returns, and larger budgets increase both the chance of successful conjecturing and the concentration of mass on high-utility conjectures.

\subsection{Query Model.}

The query model is a stateful predictive module used by each agent to estimate how its prover process is expected to perform on a target sequent under its current local library state. In our implementation, the query model does \emph{not} estimate semantic truth directly. Rather, for agent $\alpha$, it estimates:

\begin{enumerate}
    \item the probability that $\alpha$ will \emph{eventually resolve} a target sequent under its prover process, and
    \item the expected \emph{additional compute budget} required to do so, conditional on eventual success.
\end{enumerate}

Accordingly, we instantiate the query readout as a map
\[
q^\alpha : \mathcal Q^\alpha \times \mathscr L \times \mathcal S \to [0,1]\times \mathbb N,
\]
where
\[
q^\alpha(Q^\alpha_t,\mathcal L,s)
=
\big(\hat p^\alpha_t(\mathcal L,s),\hat \tau^\alpha_t(\mathcal L,s)\big).
\]
Here:
\begin{itemize}
    \item $\hat p^\alpha_t(\mathcal L,s)$ estimates the probability that agent $\alpha$'s prover process will eventually resolve $s$ from local library state $\mathcal L$,
    \item $\hat\tau^\alpha_t(\mathcal L,s)$ estimates the expected additional budget required to resolve $s$, conditional on eventual success.
\end{itemize}

\paragraph{Feature buckets.}
To make online updates statistically stable while preserving algebraic simplicity, we use a finite bucketization of query instances. For each agent $\alpha$, fix a finite bucket set $\mathcal B^\alpha$ and a deterministic bucket map
\[
b^\alpha : \mathscr L\times \mathcal S \to \mathcal B^\alpha.
\]
The bucket map may depend on any collection of state/query features, such as:
\begin{itemize}
    \item support mass or utility features of the target,
    \item dependency-completion indicators,
    \item target difficulty proxies,
    \item cumulative prior budget spent on the target,
    \item local graph-topological features,
    \item other agent-specific or state-specific summary statistics.
\end{itemize}
The framework does not prescribe the exact feature design.

\paragraph{Discrete-time survival model.}
Fix a finite horizon $H\in\mathbb N^+$. For each agent $\alpha$ and bucket $b\in\mathcal B^\alpha$, we model a discrete-time hazard sequence
\[
\lambda^\alpha_b(k)\in[0,1], \qquad k=1,\dots,H,
\]
where
\[
\lambda^\alpha_b(k)
:=
\Pr(T^\alpha_s = k \mid T^\alpha_s \ge k,\ b^\alpha(\mathcal L,s)=b).
\]
Here $T^\alpha_s$ denotes the random additional budget required for agent $\alpha$ to resolve $s$.

Equivalently, conditional on bucket $b$, the discrete survival function is
\[
S^\alpha_b(k)
:=
\Pr(T^\alpha_s > k \mid b)
=
\prod_{j=1}^{k}\big(1-\lambda^\alpha_b(j)\big),
\qquad S^\alpha_b(0):=1.
\]
Thus the probability of solving exactly at budget step $k$ is
\[
\Pr(T^\alpha_s = k \mid b)
=
\lambda^\alpha_b(k)\prod_{j=1}^{k-1}(1-\lambda^\alpha_b(j)).
\]

\paragraph{Query readout.}
Given current state $(Q^\alpha_t,\mathcal L,s)$, let
\[
b=b^\alpha(\mathcal L,s).
\]
Then the query model returns:
\[
\hat p^\alpha_t(\mathcal L,s)
=
\Pr(T^\alpha_s \le H \mid b)
=
1-S^\alpha_b(H)
=
1-\prod_{j=1}^{H}(1-\lambda^\alpha_b(j)),
\]
which is the estimated probability of eventual solve success within the modeling horizon $H$.

For the conditional expected additional budget, define
\[
\hat\tau^\alpha_t(\mathcal L,s)
=
\mathbb E[T^\alpha_s \mid T^\alpha_s \le H,\ b]
=
\frac{\sum_{k=1}^{H} k\,\Pr(T^\alpha_s = k \mid b)}
{\hat p^\alpha_t(\mathcal L,s)},
\]
whenever $\hat p^\alpha_t(\mathcal L,s)>0$. If $\hat p^\alpha_t(\mathcal L,s)=0$, we may return a fixed sentinel value (for instance $H$ or $+\infty$ in an abstract implementation).

\paragraph{Query model state.}
For each agent $\alpha$, the query model state
\[
Q^\alpha_t \in \mathcal Q^\alpha
\]
stores, for every bucket $b\in\mathcal B^\alpha$ and time index $k\in\{1,\dots,H\}$, two nonnegative sufficient statistics:
\[
N^\alpha_{t}(b,k),\qquad D^\alpha_{t}(b,k).
\]
These are interpreted as:
\begin{itemize}
    \item $D^\alpha_t(b,k)$ = the number of observed proof-attempt trajectories in bucket $b$ that remained ``at risk'' at budget step $k$,
    \item $N^\alpha_t(b,k)$ = the number of those trajectories that terminated successfully exactly at budget step $k$.
\end{itemize}

\paragraph{MAP-smoothed hazard estimates.}
To ensure numerical stability and prevent degenerate zero/one hazards early in training, we introduce fixed hyperparameters
\[
a^\alpha_{b,k}>0,\qquad c^\alpha_{b,k}>0
\]
for each bucket-time pair. These act as pseudocounts. We then define the hazard estimate by the smoothed ratio
\[
\lambda^\alpha_b(k)
=
\frac{N^\alpha_t(b,k)+a^\alpha_{b,k}}
{D^\alpha_t(b,k)+a^\alpha_{b,k}+c^\alpha_{b,k}}.
\]
This may be interpreted as the coordinatewise MAP estimate of a Bernoulli hazard with Beta prior.

\paragraph{Observed data and censoring.}
Each proof attempt produces an observation datum for the querying agent. We model a datum as
\[
d=(b,\tau,\xi)\in \mathcal B^\alpha \times \mathbb N^+ \times \{0,1\},
\]
where:
\begin{itemize}
    \item $b$ is the feature bucket of the queried instance at attempt time,
    \item $\tau$ is the additional budget consumed in that attempt,
    \item $\xi=1$ indicates that the attempt succeeded exactly at budget $\tau$,
    \item $\xi=0$ indicates that the attempt was censored (timed out / did not solve) at budget $\tau$.
\end{itemize}

Such a datum contributes additive increments to the sufficient statistics:
\[
D^\alpha_t(b,k) \leftarrow D^\alpha_t(b,k)+1
\qquad\text{for all }1\le k\le \min\{\tau,H\},
\]
and, if $\xi=1$ and $\tau\le H$, additionally
\[
N^\alpha_t(b,\tau) \leftarrow N^\alpha_t(b,\tau)+1.
\]
Hence unsuccessful attempts contribute exposure only, while successful attempts contribute both exposure and one event at the realized solve time.

\paragraph{Update rule.}
Define the update operator
\[
\Delta^\alpha : \mathcal Q^\alpha \times \mathcal D(\mathcal Q^\alpha)\to \mathcal Q^\alpha
\]
to be the additive update of these sufficient statistics induced by a single datum $d$ as above. Because the update is coordinatewise addition on count tables, it is exactly associative and commutative. Therefore, if multiple observations are received in a single timestep, the resulting batch update is well-defined independently of order.

\subsection{Market Mechanisms}

We now specify the family of concrete market mechanisms tested in this work. Each mechanism is implemented as an instantiation of the abstract market mechanism interface
\[
\mathscr M=\Big(\mathcal N,\mathfrak M,\beta,\mathcal W,\{\mathcal A^\alpha_{\mathrm{mkt}}\}_{\alpha\in\mathcal N},
\{\mathrm{Adm}^\alpha_{\mathrm{mkt}}\}_{\alpha\in\mathcal N},\mathcal T_{\mathrm{mkt}}\Big),
\]
but we restrict attention to a common implementation template so that different mechanisms may be compared under the same underlying theorem-proving environment.

\paragraph{Common portfolio structure.}
For all tested mechanisms, each agent $\alpha$ maintains a portfolio
\[
\mathcal P^\alpha_t=(c^\alpha_t,H^\alpha_t,O^\alpha_t),
\]
where:
\begin{itemize}
    \item $c^\alpha_t\in\mathbb R_{\ge 0}$ is current cash,
    \item $H^\alpha_t$ is the multiset of currently held assets / claims,
    \item $O^\alpha_t$ is the multiset of currently posted market offers.
\end{itemize}
The public market ledger $\mathcal{ML}_t$ stores the union of all currently active public offers together with any mechanism-specific public state (for example AMM inventory variables or posted bounty books).

\paragraph{Worst-case collateralization.}
To prevent agents from taking positions they cannot honor, we associate to each portfolio a worst-case marked balance
\[
\underline B^\alpha_t := \underline\beta(\mathcal P^\alpha_t)\in\mathbb R,
\]
defined as the cash position minus the largest possible future liability of the held and posted claims under the mechanism's settlement rule. We require all admissible market actions to preserve
\[
\underline B^\alpha_t \ge 0.
\]
Thus agents are fully collateralized against worst-case future settlement.

\paragraph{Cash conservation.}
All tested mechanisms satisfy internal money conservation at trade time: transferring an asset from one agent to another is accompanied by an equal and opposite transfer of cash, except where the mechanism explicitly introduces escrowed collateral or deterministic settlement transfers already implied by outstanding liabilities.

\paragraph{Public settlement events.}
All tested mechanisms settle only against publicly observable events, namely changes in the public library $\mathcal L^0_t$. In particular, contract payouts may depend only on whether a target formula or sequent has been publicly resolved by a specified time.


\paragraph{Resolution events.}
For a public target sequent $s$ and deadline $T\in\mathbb N$, define the binary event
\[
E_{s,T}(w_{0:\infty}) := \mathbf 1\big[\exists\, t\le T \text{ such that } s\in \mathcal{RS}^0_t\big].
\]
Equivalently, $E_{s,T}=1$ iff $s$ is publicly resolved by time $T$.
When working in the implemented formula-based interface, we identify a formula $\phi$ with its canonical public sequent
\[
s^0_t(\phi):=[\mathcal R^0_t\proves \phi]
\]
and write $E_{\phi,T}$ by abuse of notation.



\paragraph{Mechanism: Collateralized Bilateral Contract Market}
As an initial work in the field of MARL on VAMP's, we focus on only one particular market mechanism and study its dynamics. We do not claim that such a market mechanism is globally optimal, but rather focus on it as a simple and robust communication substrate in such verified domains such as formal mathematics.

Namely, our market mechanism is a posted-offer market over binary contingent contracts. A \emph{contract type} is a triple
\[
\chi=(s,T,\ell)\in \mathcal{RS}^{0,\mathrm{cand}}\times \mathbb N \times [0,1],
\]
where:
\begin{itemize}
    \item $s$ is a public target sequent,
    \item $T$ is a strike / deadline time,
    \item $\ell\in[0,1]$ is the maximum loss parameter.
\end{itemize}

Intuitively, $\chi=(s,T,\ell)$ is a normalized unit-notional claim on the event that $s$ becomes publicly resolved by time $T$.

\paragraph{Long and short claims.}
Each contract type $\chi$ induces two complementary positions:
\[
\chi^{+} \qquad\text{and}\qquad \chi^{-},
\]
called the long and short side, respectively. Their terminal cash payoffs are:
\[
\Pi(\chi^+) =
\begin{cases}
1-\ell, & E_{s,T}=1,\\
-\ell, & E_{s,T}=0,
\end{cases}
\qquad
\Pi(\chi^-) = -\Pi(\chi^+).
\]
Thus the long side gains when the target resolves by the deadline, while the short side gains otherwise. The magnitude of the downside of the long position is exactly $\ell$, and the upside is $1-\ell$.

\paragraph{Unit notional and quantities.}
Contracts are traded in nonnegative integer quantities. A position of quantity $q\in\mathbb N$ in $\chi^+$ has payoff $q\,\Pi(\chi^+)$, and similarly for $\chi^-$. Hence larger notional exposure is represented by larger quantities rather than by changing the contract definition itself.

\paragraph{Asset creation.}
A new quantity-$q$ contract issue of type $\chi$ is created only as a matched pair:
\[
q\chi^+ \quad\text{and}\quad q\chi^-.
\]
This pair may initially be placed into the issuing agent's holdings or posted offers. Since the long and short payoffs sum to zero statewise, pair creation introduces no net wealth.

\paragraph{Public offers.}
An offer is a tuple
\[
o=(\chi,\sigma,q,p),
\]
where $\chi$ is a contract type, $\sigma\in\{+,-\}$ is the side being offered, $q\in\mathbb N^+$ is quantity, and $p\in[0,1]$ is the per-unit transaction price. Posted offers are stored in the public ledger $\mathcal{ML}_t$.

\paragraph{Trades.}
When agent $\beta$ accepts quantity $q'\le q$ from an offer $o=(\chi,\sigma,q,p)$ posted by agent $\alpha$, the ledger removes $q'\chi^\sigma$ from $\alpha$'s posted offers and transfers it to $\beta$'s holdings, while cash is transferred in the opposite direction:
\[
c^\beta \leftarrow c^\beta - q'p,\qquad c^\alpha \leftarrow c^\alpha + q'p.
\]
The residual quantity $q-q'$ remains posted if positive.

\paragraph{Collateral and admissibility.}
For a held or posted position $x$, define its worst-case liability by
\[
L(x):=\sup_{\omega}(-\Pi_\omega(x))_+.
\]
For a portfolio $\mathcal P^\alpha=(c^\alpha,H^\alpha,O^\alpha)$, define the minimum marked balance
\[
\underline B^\alpha
:=
c^\alpha - \sum_{x\in H^\alpha\cup O^\alpha}L(x).
\]
A market action is admissible only if the resulting portfolio satisfies
\[
\underline B^\alpha \ge 0.
\]
Hence every agent is fully collateralized against the worst-case settlement of all currently held and posted obligations.

\paragraph{Settlement.}
At each timestep, any contract type $\chi=(s,T,\ell)$ with deadline $T=t$ is settled against the public event $E_{s,T}$. For each agent, the corresponding cash payoff is applied to all held and posted quantities of $\chi^\pm$, and the settled contracts are removed from portfolios and the public ledger.
